\documentclass[11pt]{article}
\usepackage[hmargin=1.05in,vmargin=0.8in]{geometry}
\usepackage{amsmath,amssymb}
\usepackage[hidelinks]{hyperref}

\title{A counterexample to Conjecture 00000007845}
\author{Feiyu (Independent researcher)}
\date{}

\begin{document}
\maketitle

\begin{abstract}
For the consistent system $I_2x=0$, one randomized Kaczmarz step leaves
an expected squared relative error of $1/2$. The Euclidean condition
number of $I_2$ is $1$, so the proposed worst-case expression evaluates
to zero.
\end{abstract}

\section*{The rate claim}
The first clause of the \href{https://github.com/The-Last-Math-Competition/The-Last-Math-Competition/blob/main/conjectures/00000007845.md}{English and Chinese statements}
asserts the exact worst-case rate $1-\kappa^{-2}$ for full-row-rank
matrices, using the expected ratio of squared errors. In particular,
for each condition number $\kappa$ this expression must bound the
expected ratio for every admissible system and nonzero current error.
A single consistent square system suffices to test that bound.

For a real row $a_i$ of $A$, let $e=x-x_*$ be the current error for a
solution $x_*$ of $Ax=b$. Projection onto the hyperplane
$a_i\cdot x=b_i$ gives the new error
\[
   e^{(i)}=e-\frac{a_i\cdot e}{\|a_i\|_2^2}a_i.
\]
The expected ratio is $\sum_i p_i\|e^{(i)}\|_2^2/\|e\|_2^2$.
The row spectral norm is its Euclidean norm. Both sampling proportional
to that norm, as in the statement, and the usual squared-row-norm
sampling give the same probabilities for rows of equal norm.

\section*{Counterexample}
Take $A=I_2$, $b=0$, $x_*=0$, and $x=(1,0)$. The rows are $(1,0)$ and
$(0,1)$, both of norm one, so each has probability $1/2$. Choosing the
first row gives $e^{(0)}=(0,0)$. Choosing the second leaves
$e^{(1)}=(1,0)$. Since $\|e\|_2^2=1$, the rate is
\[
  \tfrac12\cdot 0+\tfrac12\cdot 1=\tfrac12.
\]
The matrix has rank two and
$\kappa_2(I_2)=\|I_2\|_2\|I_2^{-1}\|_2=1$.
Thus the actual rate $1/2$ exceeds $1-\kappa_2(I_2)^{-2}=0$.
This refutes the conjunction in the conjecture independently of the
additional claim about the expected logarithmic ratio.

\section*{Lean verification}
The Lean project uses Mathlib's real matrices, matrix rank, matrix
inverse, dot product, and Euclidean operator norm. It defines the
Kaczmarz error projection and computes the expectation as a sum over
the two rows. The theorem \texttt{kaczmarz\_counterexample} produces
a full-row-rank matrix and nonzero error with a strict violation.
The final theorem \texttt{conjecture\_00000007845\_false} refutes the
stated bound for each of the two sampling rules.

Run \texttt{lake build} in \texttt{lean/}. The pinned Lean and Mathlib
versions are both v4.33.1. The final theorem's axiom dependencies are
\texttt{propext}, \texttt{Classical.choice}, and \texttt{Quot.sound}.
\end{document}
